\section*{Hoofdstuk \ref{ch:b2:measurement-statistics} --- Statistiek van chemische metingen}

\begin{solution}{exo:b2:measurement-statistics:1}
$\bar V = 62.41/5 = \qty{12.482}{mL}$; $\sum(V_i - \bar V)^2 = 0.00328$, $s = \sqrt{0.00328/4} \approx
\qty{0.029}{mL}$; $s/\sqrt5 \approx \qty{0.013}{mL}$.
\end{solution}

\begin{solution}{exo:b2:measurement-statistics:2}
$t = 2.776$ voor 4 vrijheidsgraden: $12.482 \pm 2.776 \times 0.0128$, dus $(12.48 \pm 0.04)$~mL.
\end{solution}

\begin{solution}{exo:b2:measurement-statistics:3}
De relatieve onzekerheden van de massa en van het volume zijn 0.04~\% en 0.06~\%; kwadratisch
gecombineerd geeft dat
\[\frac{u(c)}{c} = \sqrt{\Bigl(\frac{0.0002}{0.5000}\Bigr)^2 + \Bigl(\frac{0.15}{250.0}\Bigr)^2}
\approx 0.07~\%.\]
\end{solution}

\begin{solution}{exo:b2:measurement-statistics:4}
$\bar x = 1.5$, $\bar y = 0.31$, $S_{xx} = 5$, $S_{xy} = 0.99$: $b = 0.198$, $a = 0.31 - 0.198 \times 1.5
= 0.013$.
\end{solution}

\begin{solution}{exo:b2:measurement-statistics:5}
$|10.12 - 10.00|/(0.08/\sqrt5) = 0.12/0.036 \approx 3.4 > 2.776$: het verschil is significant, de
methode is vertekend (met ongeveer \qty{+0.12}{mg/L}).
\end{solution}

\begin{solution}{exo:b2:measurement-statistics:6}
$\mathrm{pH} = -\log_{10}c$, $\mathrm d\,\mathrm{pH}/\mathrm dc = -1/(c\ln10)$:
$u(\mathrm{pH}) = (u(c)/c)/\ln10 = 0.02/2.303 \approx 0.009$.
\end{solution}

\begin{solution}{exo:b2:measurement-statistics:7}
$x_{\mathrm{LOD}} = 3 \times 0.0012/0.0098 \approx \qty{0.37}{\micro g/L}$; $x_{\mathrm{LOQ}} = 10
\times 0.0012/0.0098 \approx \qty{1.2}{\micro g/L}$.
\end{solution}

\begin{solution}{exo:b2:measurement-statistics:8}
De rechte geeft \qty{3.11}{mg/L} in de gemeten oplossingen; het monster was vijfmaal verdund:
$3.11 \times 50.0/10.0 \approx \qty{15.6}{mg/L}$.
\end{solution}

\begin{solution}{exo:b2:measurement-statistics:9}
\qty{0.05}{mg/L} is de \omterm{def:b2:measurement-statistics:validation}{herhaalbaarheid}, \qty{0.15}{mg/L} de \omterm{def:b2:measurement-statistics:validation}{reproduceerbaarheid}. Elk laboratorium heeft zijn
eigen kleine \omterm{def:b2:measurement-statistics:validation}{bias} (ijking, standaarden, instrument); die biassen verschillen tussen laboratoria en tellen op
bij de spreiding, dus de \omterm{def:b2:measurement-statistics:validation}{reproduceerbaarheid} is altijd de grootste.
\end{solution}

\begin{solution}{exo:b2:measurement-statistics:10}
$\partial S/\partial a = 0$ luidt $\sum(y_i - a - bx_i) = 0$: de som van de \omterm{def:b2:measurement-statistics:calibration}{residuen} is nul. Delen door
$n$ geeft $\bar y = a + b\bar x$, dus $(\bar x, \bar y)$ ligt op de rechte.
\end{solution}

\begin{solution}{exo:b2:measurement-statistics:11}
Eén stap:
\[\sqrt{(0.006/1.000)^2 + (0.08/100.0)^2} \approx 0.61~\%.\]
Eén stap van tienmaal:
\[\sqrt{(0.02/10.00)^2 + (0.08/100.0)^2} \approx 0.22~\%,\]
en twee daarvan kwadratisch gecombineerd geven $0.22\sqrt2 \approx 0.30~\%$. De route in twee stappen is twee keer zo precies: de kleine pipet domineert de route
in één stap.
\end{solution}

\begin{solution}{exo:b2:measurement-statistics:12}
$z = 1.6/\sqrt{0.4^2 + 0.5^2} \approx 2.5 > 2$: niet verenigbaar; (minstens) één laboratorium heeft een \omterm{def:b2:measurement-statistics:validation}{bias}.
Met uitgebreide onzekerheden ($k = 2$) geeft het eerste $9.6 \pm 0.8$, wat 10 bevat; het tweede
$11.2 \pm 1.0$, volledig boven 10. Voor er een oordeel over het water komt, moet de onenigheid opgelost worden,
bijvoorbeeld met een referentiemateriaal.
% ledger: who:lead
\end{solution}

\begin{solution}{pb:b2:measurement-statistics:1}
\textbf{1.} Het zuur verandert de atomisering en de respons; standaarden en monsters moeten dezelfde
matrix hebben opdat de helling geldt.
\textbf{2.} $\bar x = 10$, $\bar y = 0.1004$, $S_{xx} = 250$, $S_{xy} = 2.445$: $b = 0.00978$ per
\unit{\micro g/L}, $a = 0.1004 - 0.0978 = 0.0026$.
\textbf{3.} $+0.0004$, $-0.0005$, $+0.0006$, $-0.0013$, $+0.0008$.
\textbf{4.} De tekens wisselen af, geen kromming: de rechte volstaat over 0--\qty{20}{\micro g/L}.
\textbf{5.} $s_{y/x} = \sqrt{3.1 \times 10^{-6}/3} \approx 0.00102$.
\textbf{6.} De blanco (zuur, water, oven) geeft zelf een kleine absorbantie.
\textbf{7.} De gevoeligheid: absorbantie per \unit{\micro g/L} lood.
\textbf{8.} $\bar y_0 = 0.1010$; $s = 0.0030$.
\textbf{9.} $x_0 = (0.1010 - 0.0026)/0.00978 \approx \qty{10.06}{\micro g/L}$.
\textbf{10.} $s_{x_0} = (0.00102/0.00978)\sqrt{1/3 + 1/5 + (0.1010 - 0.1004)^2/(0.00978^2 \times 250)}
\approx 0.104 \times 0.730 \approx \qty{0.076}{\micro g/L}$.
\textbf{11.} $n - 2 = 3$; $t = 3.182$.
\textbf{12.} $10.06 \pm 3.182 \times 0.076 = (10.06 \pm 0.24)$~\unit{\micro g/L}.
\textbf{13.} De term $1/m$ daalt van 1 naar $1/3$: $s_{x_0}$ daalt met ongeveer een derde.
\textbf{14.} $f = 50.50/50.0 = 1.010$; $10.06 \times 1.010 \approx \qty{10.16}{\micro g/L}$.
\textbf{15.} $f = 1 + V_2/V_1$: $u^2(f) = (u(V_2)/V_1)^2 + (V_2u(V_1)/V_1^2)^2 = (0.005/50.0)^2 +
(0.50 \times 0.03/2500)^2$, $u(f) \approx 1.0 \times 10^{-4}$, relatief 0.01~\%.
\textbf{16.} Nee: 0.01~\% tegenover $0.076/10.06 \approx 0.75~\%$.
\textbf{17.} $s_{\mathrm{blank}} \approx 0.00115$; $x_{\mathrm{LOD}} = 3 \times 0.00115/0.00978 \approx
\qty{0.35}{\micro g/L}$.
\textbf{18.} $10 \times 0.00115/0.00978 \approx \qty{1.2}{\micro g/L}$.
\textbf{19.} Ja, ruimschoots.
\textbf{20.} Ja: van 9.92 tot \qty{10.40}{\micro g/L} in het kraanwater.
\textbf{21.} $(0.104 - 0.0026)/0.00978 \times 1.010 \approx \qty{10.5}{\micro g/L}$, en, op zich gelezen,
een zelfverzekerd ``boven de richtwaarde'' dat de gegevens niet ondersteunen.
\textbf{22.} De regel waarmee het interval opgebouwd wordt, vangt de ware waarde in 95~\% van de reeksen waarop
hij toegepast wordt; het is geen kans over deze ene waarde.
\textbf{23.} Meer herhaalde metingen en meer standaarden rond \qty{10}{\micro g/L} (de termen $1/m$ en
$1/n$), of een gevoeligere methode; het interval krimpt maar langzaam, als $1/\sqrt{\text{aantal}}$.
\textbf{24.} Door een gecertificeerd referentiewater te analyseren, of door een bekende hoeveelheid lood aan het monster
toe te voegen en de terugvinding te controleren.
\textbf{25.} Rond \qty{10}{\micro g/L} bedraagt de onzekerheid van routinemethoden enkele procenten, zoals hier: een
lagere richtwaarde zou door gewone laboratoria niet betrouwbaar gecontroleerd kunnen worden.
\textbf{26.} $\boldsymbol{(10.16 \pm 0.24)}$~\unit{\micro g/L} bij 95~\%: het interval bevat
\qty{10}{\micro g/L}, dus de meting toont noch aan dat het water de richtwaarde overschrijdt, noch dat het er met een
marge aan voldoet.
\end{solution}
